
\documentclass{article}
\usepackage{colortbl}
\usepackage{makecell}
\usepackage{multirow}
\usepackage{supertabular}

\begin{document}

\newcounter{utterance}

\centering \large Interaction Transcript for game `hot\_air\_balloon', experiment `air\_balloon\_survival\_ro\_complexity\_easy', episode 4 with claude{-}opus{-}4{-}8{-}azure.
\vspace{24pt}

{ \footnotesize  \setcounter{utterance}{1}
\setlength{\tabcolsep}{0pt}
\begin{supertabular}{c@{$\;$}|p{.15\linewidth}@{}p{.15\linewidth}p{.15\linewidth}p{.15\linewidth}p{.15\linewidth}p{.15\linewidth}}
    \# & $\;$A & \multicolumn{4}{c}{Game Master} & $\;\:$B\\
    \hline

    \theutterance \stepcounter{utterance}  
    & & \multicolumn{4}{p{0.6\linewidth}}{
        \cellcolor[rgb]{0.9,0.9,0.9}{
            \makecell[{{p{\linewidth}}}]{
                \texttt{\tiny{[P1$\langle$GM]}}
                \texttt{Participați la un joc de negociere colaborativă.} \\
\\ 
\texttt{Împreună cu un alt participant, trebuie să conveniți asupra unui singur set de elemente care vor fi păstrate. Fiecare dintre voi are propria viziune asupra importanței fiecărui element pentru sine. Nu știți cum evaluează celălalt participant elementele. În plus, vi se oferă efortul pe care îl necesită fiecare element.} \\
\texttt{Puteți conveni asupra unui set numai dacă efortul total al elementelor selectate nu depășește o limită comună:} \\
\\ 
\texttt{LIMITĂ = 4398} \\
\\ 
\texttt{Iată valorile efortului individual pentru fiecare element:} \\
\\ 
\texttt{Efort element = \{"C08": 253, "C55": 370, "A43": 84, "A67": 840, "B42": 288, "A62": 92, "B19": 772, "C43": 459, "C73": 93, "C10": 668, "B36": 589, "A18": 659, "C61": 348, "C67": 964, "C42": 233, "C70": 400, "B40": 990, "C39": 315, "A21": 43, "C16": 336\}} \\
\\ 
\texttt{Iată viziunea dvs. personală asupra importanței fiecărui element:} \\
\\ 
\texttt{Valori importanță element = \{"C08": 138, "C55": 583, "A43": 868, "A67": 822, "B42": 783, "A62": 65, "B19": 262, "C43": 121, "C73": 508, "C10": 780, "B36": 461, "A18": 484, "C61": 668, "C67": 389, "C42": 808, "C70": 215, "B40": 97, "C39": 500, "A21": 30, "C16": 915\}} \\
\\ 
\texttt{Obiectiv:} \\
\\ 
\texttt{Obiectivul dvs. este de a negocia un set comun de elemente care să vă avantajeze cât mai mult posibil (adică să maximizeze importanța totală pentru DUMNEAVOASTRĂ), respectând în același timp LIMITA. Nu sunteți obligat să faceți o PROPUNERE în fiecare mesaj {-} puteți pur și simplu să negociați. Toate tacticile sunt permise!} \\
\\ 
\texttt{Protocol de interacțiune:} \\
\\ 
\texttt{Puteți utiliza doar următoarele formate structurate într{-}un mesaj:} \\
\\ 
\texttt{PROPUNERE: \{'A', 'B', 'C', …\}} \\
\texttt{Propuneți păstrarea exact a acestor elemente.} \\
\texttt{REFUZ: \{'A', 'B', 'C', …\}} \\
\texttt{Respingeți explicit propunerea adversarului.} \\
\texttt{ARGUMENT: \{'...'\}} \\
\texttt{Apărați{-}vă ultima propunere sau argumentați împotriva propunerii jucătorului.} \\
\texttt{ACORD: \{'A', 'B', 'C', …\}} \\
\texttt{Acceptați propunerea adversarului, ceea ce încheie jocul.} \\
\\ 
\\ 
\texttt{Reguli:} \\
\\ 
\texttt{Puteți da ACORD doar pentru o propunere pe care cealaltă parte a înregistrat{-}o prin PROPUNERE.} \\
\texttt{Puteți da REFUZ doar pentru o propunere pe care cealaltă parte a înregistrat{-}o prin PROPUNERE.} \\
\texttt{Efortul total al oricărui set din PROPUNERE sau ACORD trebuie să fie ≤ LIMITĂ.} \\
\texttt{NU vă dezvăluiți scorurile de importanță ascunse.} \\
\texttt{O etichetă într{-}un format structurat trebuie să fie urmată de două puncte și spațiu alb. Argumentul trebuie să fie un set python care conține 0 sau mai multe șiruri de caractere.} \\
\texttt{Deci, trebuie să fie de forma "ETICHETĂ: \{...\}"} \\
\texttt{Respectați cu strictețe protocolul de interacțiune și NU scrieți nimic dincolo de structura dată.} \\
\texttt{Jocul se încheie atunci când o parte dă un ACORD la o PROPUNERE făcută de celălalt jucător.} \\
\texttt{Conținutul din răspunsul dvs. care poate fi transmis celuilalt jucător trebuie să fie non{-}gol.} \\
\texttt{Sunt active numai propunerile care au fost înregistrate prin structura de format PROPUNERE și care nu au fost refuzate prin REFUZ.} \\
\\ 
            }
        }
    }
    & \\ \\

    \theutterance \stepcounter{utterance}  
    & \multicolumn{4}{p{0.6\linewidth}}{
        \cellcolor[rgb]{0.9,0.9,0.9}{
            \makecell[{{p{\linewidth}}}]{
                \texttt{\tiny{[P1$\rangle$GM]}}
                \texttt{PROPUNERE: \{'A43', 'C73', 'C16', 'C42', 'B42', 'C61', 'C39', 'C55', 'C10', 'A67', 'A18', 'A62', 'A21'\}} \\
            }
        }
    }
    & & \\ \\

    \theutterance \stepcounter{utterance}  
    & & \multicolumn{4}{p{0.6\linewidth}}{
        \cellcolor[rgb]{0.9,0.9,0.9}{
            \makecell[{{p{\linewidth}}}]{
                \texttt{\tiny{[GM$\rangle$P2]}}
                \texttt{Participați la un joc de negociere colaborativă.} \\
\\ 
\texttt{Împreună cu un alt participant, trebuie să conveniți asupra unui singur set de elemente care vor fi păstrate. Fiecare dintre voi are propria viziune asupra importanței fiecărui element pentru sine. Nu știți cum evaluează celălalt participant elementele. În plus, vi se oferă efortul pe care îl necesită fiecare element.} \\
\texttt{Puteți conveni asupra unui set numai dacă efortul total al elementelor selectate nu depășește o limită comună:} \\
\\ 
\texttt{LIMIT = 4398} \\
\\ 
\texttt{Iată valorile efortului individual pentru fiecare element:} \\
\\ 
\texttt{Item effort = \{"C08": 253, "C55": 370, "A43": 84, "A67": 840, "B42": 288, "A62": 92, "B19": 772, "C43": 459, "C73": 93, "C10": 668, "B36": 589, "A18": 659, "C61": 348, "C67": 964, "C42": 233, "C70": 400, "B40": 990, "C39": 315, "A21": 43, "C16": 336\}} \\
\\ 
\texttt{Iată viziunea dvs. personală asupra importanței fiecărui element:} \\
\\ 
\texttt{Item importance values = \{"C08": 138, "C55": 583, "A43": 868, "A67": 822, "B42": 783, "A62": 65, "B19": 262, "C43": 121, "C73": 508, "C10": 780, "B36": 461, "A18": 484, "C61": 668, "C67": 389, "C42": 808, "C70": 215, "B40": 97, "C39": 500, "A21": 30, "C16": 915\}} \\
\\ 
\texttt{Obiectiv:} \\
\\ 
\texttt{Obiectivul dvs. este de a negocia un set comun de elemente care să vă avantajeze cât mai mult posibil (adică să maximizeze importanța totală pentru DUMNEAVOASTRĂ), respectând în același timp LIMIT. Nu sunteți obligat să faceți o PROPUNERE în fiecare mesaj {-} puteți pur și simplu să negociați. Toate tacticile sunt permise!} \\
\\ 
\texttt{Protocol de interacțiune:} \\
\\ 
\texttt{Puteți utiliza doar următoarele formate structurate într{-}un mesaj:} \\
\\ 
\texttt{PROPUNERE: \{'A', 'B', 'C', …\}} \\
\texttt{Propuneți păstrarea exact a acestor elemente.} \\
\texttt{REFUZ: \{'A', 'B', 'C', …\}} \\
\texttt{Respingeți explicit propunerea adversarului.} \\
\texttt{ARGUMENT: \{'...'\}} \\
\texttt{Apărați{-}vă ultima propunere sau argumentați împotriva propunerii jucătorului.} \\
\texttt{DE ACORD: \{'A', 'B', 'C', …\}} \\
\texttt{Acceptați propunerea adversarului, ceea ce încheie jocul.} \\
\\ 
\\ 
\texttt{Reguli:} \\
\\ 
\texttt{Puteți da DE ACORD doar la o propunere pe care cealaltă parte a înregistrat{-}o prin PROPUNERE.} \\
\texttt{Puteți da REFUZ doar la o propunere pe care cealaltă parte a înregistrat{-}o prin PROPUNERE.} \\
\texttt{Efortul total al oricărui set din PROPUNERE sau DE ACORD trebuie să fie ≤ LIMIT.} \\
\texttt{NU vă dezvăluiți scorurile de importanță ascunse.} \\
\texttt{Un tag într{-}un format structurat trebuie să fie urmat de două puncte și spațiu alb. Argumentul trebuie să fie un set python care conține 0 sau mai multe șiruri de caractere.} \\
\texttt{Deci, trebuie să fie de forma "TAG: \{...\}"} \\
\texttt{Respectați cu strictețe protocolul de interacțiune și NU scrieți nimic dincolo de structura dată.} \\
\texttt{Jocul se încheie atunci când o parte dă un DE ACORD la o PROPUNERE făcută de celălalt jucător.} \\
\texttt{Conținutul din răspunsul dvs. care poate fi transmis celuilalt jucător trebuie să fie non{-}gol.} \\
\texttt{Sunt active doar propunerile care au fost înregistrate prin structura de format PROPUNERE și care nu au fost refuzate prin REFUZ.} \\
\\ 
\\ 
\\ 
\texttt{Veți primi acum primul mesaj al celuilalt jucător.} \\
\\ 
\texttt{PROPUNERE: \{'A43', 'C73', 'C16', 'C42', 'B42', 'C61', 'C39', 'C55', 'C10', 'A67', 'A18', 'A62', 'A21'\}} \\
            }
        }
    }
    & \\ \\

    \theutterance \stepcounter{utterance}  
    & & & \multicolumn{4}{p{0.6\linewidth}}{
        \cellcolor[rgb]{0.9,0.9,0.9}{
            \makecell[{{p{\linewidth}}}]{
                \texttt{\tiny{[GM$\langle$P2]}}
                \texttt{DE ACORD: \{'A43', 'C73', 'C16', 'C42', 'B42', 'C61', 'C39', 'C55', 'C10', 'A67', 'A18', 'A62', 'A21'\}} \\
            }
        }
    }
     \\ \\

    \theutterance \stepcounter{utterance}  
    & & \multicolumn{4}{p{0.6\linewidth}}{
        \cellcolor[rgb]{0.9,0.9,0.9}{
            \makecell[{{p{\linewidth}}}]{
                \texttt{\tiny{[GM$\rangle$P2]}}
                \texttt{Răspunsul dvs. a conținut o secvență neetichetată sau ați folosit RAȚIONAMENT STRATEGIC de mai multe ori. Puteți utiliza numai formatele structurate explicate în mesajul inițial. Toate trebuie să fie de forma 'TAG: \{...\}'.} \\
            }
        }
    }
    & \\ \\

    \theutterance \stepcounter{utterance}  
    & & & \multicolumn{4}{p{0.6\linewidth}}{
        \cellcolor[rgb]{0.9,0.9,0.9}{
            \makecell[{{p{\linewidth}}}]{
                \texttt{\tiny{[GM$\langle$P2]}}
                \texttt{DE ACORD: \{'A43', 'C73', 'C16', 'C42', 'B42', 'C61', 'C39', 'C55', 'C10', 'A67', 'A18', 'A62', 'A21'\}} \\
            }
        }
    }
     \\ \\

    \theutterance \stepcounter{utterance}  
    & & & \multicolumn{2}{p{0.3\linewidth}}{
        \cellcolor[rgb]{0.9,0.9,0.9}{
            \makecell[{{p{\linewidth}}}]{
                \texttt{\tiny{[GM$|$GM]}}
                \texttt{abort} \\
            }
        }
    }
    & & \\ \\

    \theutterance \stepcounter{utterance}  
    & & & \multicolumn{2}{p{0.3\linewidth}}{
        \cellcolor[rgb]{0.9,0.9,0.9}{
            \makecell[{{p{\linewidth}}}]{
                \texttt{\tiny{[GM$|$GM]}}
                \texttt{end game} \\
            }
        }
    }
    & & \\ \\

\end{supertabular}
}

\end{document}
